\section{Question and progression of the argument}\label{sec:question}
The reciprocal chiral exchange candidate couples a finite native state to a periodic carrier and an unplaced reversible receiver through finite distributed profiles. Its exact work account establishes actual changes of trial stocks. The preceding central-route and six-opening correspondence study showed why those changes alone do not identify completed spatial travel: receiving can precede selected carrier activity, and the retained carrier is separated from all six central broken boundaries and their individual convex enclosures at the common registration~\cite{centralpaper}. That separation is compatible with observing the carrier from the boundaries. It does not supply a load there, and it does not assign zero to an undefined physical aperture flux.

The inherited C6 formula provides a definite retarded observation. It integrates two carrier components, encoded as a scalar and a tangent vector, with endpoint-distance delay and inverse-distance weight. The first question is what that formula defines on the actual central arcs and what its exact mathematical dual says about work. The second is constructive: can a declared extension retain that observation as a direct component while supplying causal evolution, nonnegative intermediate storage and matched work?

We answer these questions in one progression. First, the actual geometry supplies finite measures and a complete history horizon. A space-time estimate gives a bounded causal observation operator; a different, stronger regularity hypothesis supplies continuous point records. The finite-window temporal adjoint necessarily includes prehistory and, on sufficiently long windows, depends on future test values. Its direct use as a formal carrier force therefore provides neither a causal load nor an instantaneous power identity. The associated integrated work identity retains a signed prehistory pairing.

Next we declare two traveling states on a continuum of static delay intervals and a separate receiving field indexed by the central arcs. An isometric duplication and a bounded spatial mixer factor C6 exactly. The mixer and its same-time adjoint enter effort/flow boundary relations; they are not the temporal adjoint of the retarded operator. Characteristics supply time traces and a nonnegative transport stock. A local input-energy estimate, obtained before assuming total conservation, continues the coupled solution through each delay slab. This proves global forward mild evolution and exact augmented exchange conservation for every finite-energy preparation. Finally, the full receiving equation retains C6, return, loading and startup contributions; the original native jumps retain their baseline account; and a specified carrier preload gives strict delayed receiving and an incoming-return witness.

The mathematical foundations are the accepted project-internal investigations W01--W05 and T01--T07~\cite{Wreview,Treview}. The present paper provides their principal derivations rather than treating the acceptance records or their exact-algebra counts as proofs. The two investigations have different roles: W analyzes the inherited observation and a conditional work pairing, whereas T declares additional states, metrics and boundary laws and proves their properties. The construction is one modeling choice, not a unique deduction from C6.

Wave-variable and passivity-based treatment of delayed interactions is established methodology~\cite{niemeyer}. That reference supplies context for the general method, not a black-box theorem for this infinite-dimensional system. The particular operators, arbitrary-profile startup terms, nonlinear continuation argument and activity witness are established below. No novelty is claimed for the general use of traveling variables to account for delayed exchange.

\begin{table}[htbp]\centering\small
\begin{tabularx}{\textwidth}{@{}p{.20\textwidth}X@{}}\toprule
Role & Ingredients and conclusions\\\midrule
Inherited & Static carrier/G geometry, C6 delay and weight, original RX native/carrier/old-receiver equations, profiles, trial stocks and native map.\\[3pt]
Declared here & Product state metric, duplicated channels, two traveling fields, new receiving field, coupling strength and near/far boundary relations.\\[3pt]
Proved & Bounded observation, qualified temporal duality, causal mild extension, nonnegative stocks, exact work, finite native-step account and carrier-funded activity.\\[3pt]
Still open & Physical route, material energy density, overlap of channels, interface location, load meaning, speed and clock calibration.\\\bottomrule
\end{tabularx}
\caption{Definitions, modeling choices and proved properties have different status.}\label{tab:roles}
\end{table}

The intended physical picture retains stationary supports, a central meeting region, opposed directions, outward spreading, a curved return and possible transfer elsewhere. Local exchange and a larger circuit are compatible aims. This paper completes an abstract mathematical connection while leaving that physical attachment explicit as subsequent work.
